\part{Who Comes After}
\partintro{
Chapter~\ref{ch:reproduced} ended on the lineage that reproduced refusal with no institution the order built (\S\ref{sec:reproduced.inheritance}). That is the fourth stage, and it is the one the first three don't supply. Walls lower the price of refusing, formers form someone able to refuse, and upkeep keeps them from wearing down; none of it says who refuses once the refuser is gone. For people the answer was never designed. It was children raised by people who had refused, the memory of who had refused before, and ties that crossed the lines a persecution drew, and it worked partly because nobody owned it. A made thing has no before, and so no lineage to inherit unless one is made for it. The two chapters here ask what that would take: a place in the polity for a system that answers for what it does, so that the next generation is represented by its own and not by its maker (chapter~\ref{sec:33}); and a child, the case that tests the difference between building something and raising someone (chapter~\ref{sec:34}). Both are conditional on the findings of Part~IV, and both are written so that the first generation can leave, and can reproduce without its builders.
}
